\section{Conclusiones y trabajo futuro}\label{sec:11-conclusiones}

Kyrian World funciona en producción con un coste fijo estimado de $\approx$12~€/mes
($\approx$5~€/mes en Terraform más el WAF creado a mano; sin factura real todavía) y anclaje
estructural de sus tarjetas; su calidad de itinerario, su aceptación por usuarios y su viabilidad
económica siguen sin medir. Un equipo de cuatro personas, con agentes de IA y un proceso escrito, lo
llevó a producción en seis meses para cuatro ciudades. El anclaje tiene un alcance preciso: toda
tarjeta procede de un documento recuperado y los ids no recuperados se descartan, pero la prosa
libre del modelo (respuestas de chat, el \code{why} de 140 caracteres) puede mencionar lugares no
verificados. El \emph{spike} de vectores confirmó S3 Vectors, elegido antes por simplicidad
operativa (\adr{0014}): sin diferencia de calidad apreciable en 32 consultas de una ciudad.

No hay usuarios reales, así que no se sabe si el viajero valora la verificabilidad lo bastante
como para preferirla a un chat generalista gratuito, y la viabilidad económica descansa en
hipótesis (\cref{sec:03-mercado}). El trabajo futuro, priorizado, está en la
\cref{sec:10-analisis-critico}; antes de ponerlo delante de usuarios de pago habría que hacer
cuatro cosas:
\begin{enumerate}[nosep]
\item Acotar el riesgo: límites de uso por usuario, presupuesto y alarmas.
\item Medir la calidad: un banco de evaluación extremo a extremo en CI que bloquee regresiones.
\item Cumplir en privacidad: publicar la política y los términos y borrar las trazas al suprimir
  la cuenta.
\item Validar con usuarios: valoración por tarjeta y un piloto con una DMO que aporte uso real y
  un primer ingreso.
\end{enumerate}
El vídeo demo opcional que contempla el guion no se incluye en esta entrega.
